\chapter{Pruebas y calidad}
\label{cap:pruebas-calidad}

\section{Pruebas automatizadas del backend}

La suite completa del backend, ejecutada con \texttt{mvn clean verify}
sobre el commit exacto del tag hist\'orico e inmutable \texttt{v1.0.0}
(\texttt{0d5cd52}) y verificada directamente contra
\texttt{Backend/target/surefire-reports/} de esa reproducci\'on
(\texttt{docs/mediciones/sec/reproduccion-v1.0.0/}), reporta:

\begin{table}[H]
  \centering
  \caption{Resultado real de la suite de pruebas del backend sobre el commit del tag \texttt{v1.0.0} (evidencia hist\'orica; ver estado del \texttt{HEAD} actual m\'as abajo).}
  \label{tab:pruebas-final}
  \begin{tabular}{@{}ll@{}}
    \toprule
    M\'etrica & Valor \\
    \midrule
    Pruebas ejecutadas & 205 \\
    Fallos & 0 \\
    Errores & 0 \\
    Omitidas & 0 \\
    Resultado del build & \texttt{BUILD SUCCESS} \\
    \bottomrule
  \end{tabular}
\end{table}

La figura~\ref{fig:maven-verify-final} resume este resultado, regenerada
autom\'aticamente desde ese mismo log crudo versionado mediante
\texttt{docs/informe/scripts/generar-figuras-evidencia.py}, sin capturas
manuales.

% EVIDENCIA-JAIME-01 (regenerada automaticamente; sustituye a la figura
% historica figuras/jaime/01-maven-verify.png, que se conserva en el
% repositorio como artefacto historico pero ya no se referencia aqui)
% Ruta esperada: figuras/jaime/07-maven-verify-final.png
\evidencia{figuras/jaime/07-maven-verify-final.png}%
{Resultado real de \texttt{mvn clean verify} sobre el commit del tag \texttt{v1.0.0}: 205 pruebas, 0 fallos, 0 errores, 0 omitidas, \texttt{BUILD SUCCESS} (regenerada por \texttt{docs/informe/scripts/generar-figuras-evidencia.py} desde el log crudo versionado).}%
{fig:maven-verify-final}

Estas 205 pruebas corresponden al estado hist\'orico e inmutable del tag
\texttt{v1.0.0}. Ejecutar \texttt{mvn clean verify} sobre el
\texttt{HEAD} actual de la rama de recalificaci\'on -- que incorpora
pruebas nuevas sobre esa misma base, sin sustituirla -- produce
\textbf{217} pruebas (0 fallos, 0 errores, 0 omitidas),
\texttt{BUILD SUCCESS}; ambas cifras son correctas dentro de su alcance
temporal respectivo. Detalle de la diferencia (205~$\rightarrow$~217) y
reproducci\'on de ambos estados en
\texttt{docs/mediciones/TEST-COUNT-PROVENANCE.md} y en la tabla de
\texttt{README.md} (secci\'on ``Estado t\'ecnico verificado'').

\section{Cobertura JaCoCo}
\label{sec:jacoco-final}

JaCoCo~\cite{jacoco-docs} mide y verifica autom\'aticamente la cobertura
con una regla \texttt{BUNDLE} sobre todo el m\'odulo backend. El umbral
autom\'atico configurado en \texttt{Backend/pom.xml} se elev\'o de 60\,\%
(Tercera Entrega) a \textbf{70\,\%} para v1.0.0.

\begin{table}[H]
  \centering
  \caption{Cobertura medida por \texttt{jacoco:check} sobre el commit del tag \texttt{v1.0.0} (evidencia hist\'orica; fuente: \texttt{docs/mediciones/jacoco/METRICS.md}). Ver estado del \texttt{HEAD} actual m\'as abajo.}
  \label{tab:jacoco-final}
  \begin{tabular}{@{}lrrrr@{}}
    \toprule
    Contador & Cubierto & No cubierto & Cobertura & Umbral \\
    \midrule
    LINE & 885 & 79 & 91{,}80\,\% & 70\,\% \\
    BRANCH & 181 & 47 & 79{,}39\,\% & 70\,\% \\
    COMPLEXITY & 299 & 77 & 79{,}52\,\% & 60\,\% \\
    \bottomrule
  \end{tabular}
\end{table}

La figura~\ref{fig:jacoco-resumen-final} muestra LINE y BRANCH frente al
umbral de 70\,\% (COMPLEXITY, con umbral distinto de 60\,\%, no se
incluye), regenerada autom\'aticamente desde \texttt{jacoco.csv} (45
clases) mediante el mismo script.

% EVIDENCIA-JAIME-02 (regenerada automaticamente; sustituye a la figura
% historica figuras/jaime/02-jacoco-resumen.png, que se conserva en el
% repositorio como artefacto historico pero ya no se referencia aqui)
% Ruta esperada: figuras/jaime/08-jacoco-resumen-final.png
\evidencia{figuras/jaime/08-jacoco-resumen-final.png}%
{Cobertura JaCoCo real sobre las 45 clases del backend: LINE 91{,}80\,\% y BRANCH 79{,}39\,\%, ambas por encima del umbral de calidad de 70\,\% (regenerada por \texttt{docs/informe/scripts/generar-figuras-evidencia.py} desde \texttt{jacoco.csv}).}%
{fig:jacoco-resumen-final}

Los tres contadores superan con margen su umbral configurado. La cobertura
mide ejecuci\'on de pruebas, no ausencia de defectos: una l\'inea cubierta
por una prueba no implica que est\'e libre de errores l\'ogicos no
contemplados por la aserci\'on de esa prueba.

Esta cobertura (LINE 91{,}80\,\%, BRANCH 79{,}39\,\%) corresponde al
estado hist\'orico e inmutable del tag \texttt{v1.0.0}. Recalculada sobre
el \texttt{HEAD} actual de la rama de recalificaci\'on, \texttt{jacoco.csv}
arroja LINE~\textbf{90{,}22\,\%} (876/971) y BRANCH~\textbf{76{,}27\,\%}
(180/236); ambas siguen superando con margen el umbral de 70\,\% del
\texttt{pom.xml}. Las diferencias corresponden al estado actual del
c\'odigo y de su suite de pruebas respecto del corte hist\'orico
\texttt{v1.0.0} (detalle en \texttt{README.md}, secci\'on ``Estado
t\'ecnico verificado'').

\section{Rendimiento (k6)}
\label{sec:rendimiento-final}

Herramienta de carga: k6~\cite{k6-docs}. La metodolog\'ia sigue el mismo
patr\'on de \emph{benchmarking} t\'ecnico contra un artefacto web real
descrito por Hasselbring~\cite{hasselbring2021benchmarking} como est\'andar
emp\'irico de ingenier\'ia de software, y usa los percentiles p95/p99 como
marco de referencia para interpretar el rendimiento de aplicaciones web
de negocio, en la l\'inea propuesta por Rempel~\cite{rempel2015webperf}.
Para v1.0.0 se ampli\'o el n\'umero de corridas de 3+3 (Tercera Entrega) a
\textbf{5 en fr\'io + 5 en caliente}, contra \texttt{GET /api/mascotas}
sobre HTTPS/TLS~1.3 real, 50~VUs, \emph{ramp-up} de 5~s y carga sostenida
de 30~s por corrida.

\begin{table}[H]
  \centering
  \tiny
  \caption{Resultado real de las diez corridas de k6 (fuente: \texttt{docs/mediciones/perf/REPORT.md}).}
  \label{tab:k6-final}
  \begin{tabular}{@{}lrrrrrrr@{}}
    \toprule
    Corrida & n & Media (ms) & Mediana (ms) & p95 (ms) & p99 (ms) & Error (\%) & Throughput (req/s) \\
    \midrule
    caliente-01 & 3196 & 10{,}54 & 8{,}68 & 20{,}15 & 30{,}68 & 0{,}0 & 91{,}17 \\
    caliente-02 & 3226 & 7{,}19 & 6{,}63 & 9{,}70 & 17{,}18 & 0{,}0 & 92{,}09 \\
    caliente-03 & 3226 & 6{,}60 & 6{,}01 & 10{,}86 & 16{,}13 & 0{,}0 & 92{,}13 \\
    caliente-04 & 3236 & 5{,}21 & 4{,}92 & 6{,}83 & 10{,}12 & 0{,}0 & 92{,}40 \\
    caliente-05 & 3236 & 5{,}40 & 4{,}96 & 7{,}81 & 12{,}59 & 0{,}0 & 92{,}41 \\
    fr\'io-01 & 3206 & 10{,}51 & 8{,}47 & 21{,}10 & 33{,}14 & 0{,}0 & 91{,}49 \\
    fr\'io-02 & 3211 & 9{,}37 & 7{,}60 & 17{,}23 & 25{,}79 & 0{,}0 & 91{,}62 \\
    fr\'io-03 & 3204 & 10{,}80 & 8{,}43 & 22{,}13 & 33{,}40 & 0{,}0 & 91{,}36 \\
    fr\'io-04 & 3206 & 10{,}71 & 8{,}50 & 20{,}16 & 32{,}65 & 0{,}0 & 91{,}45 \\
    fr\'io-05 & 3199 & 10{,}61 & 7{,}88 & 20{,}71 & 40{,}87 & 0{,}0 & 91{,}51 \\
    \bottomrule
  \end{tabular}
\end{table}

La tasa de error fue \textbf{0{,}0\,\%} en las diez corridas, sin
excepci\'on. Intervalos de confianza al 95\,\% (distribuci\'on t de
Student) reportados por corrida en \texttt{docs/mediciones/perf/REPORT.md};
por ejemplo, caliente-01: $[10{,}30, 10{,}78]$~ms. Adem\'as del IC~95\,\%,
se aplic\'o la prueba de Wilcoxon pareada entre condiciones fr\'ia y
caliente (por \'indice de llegada), con tama\~no de efecto $r$ de
Rosenthal: el efecto es peque\~no en la corrida~01 ($r=-0{,}07$) y grande en
las corridas~03--04 ($r=-0{,}65$ a $-0{,}84$), consistente con una cach\'e
que reduce la latencia de forma creciente conforme el sistema se estabiliza
entre corridas consecutivas.

% EVIDENCIA-FRED-06
% Ruta esperada: figuras/fred/06-performance-report.png
\evidencia{figuras/fred/06-performance-report.png}%
{Vista del reporte agregado de rendimiento.}%
{fig:performance-report-final}

% EVIDENCIA-FRED-CALIENTE
% Ruta esperada: figuras/fred/Rendimiento-PruebaCaliente.png
\evidencia{figuras/fred/Rendimiento-PruebaCaliente.png}%
{Salida de consola de una corrida de k6 en caliente.}%
{fig:rendimiento-caliente-final}

% EVIDENCIA-FRED-FRIO
% Ruta esperada: figuras/fred/Rendimiento-PruebaFria.png
\evidencia{figuras/fred/Rendimiento-PruebaFria.png}%
{Salida de consola de una corrida de k6 en fr\'io.}%
{fig:rendimiento-frio-final}

\section{Usabilidad (SUS)}
\label{sec:sus-final}

El instrumento System Usability Scale de Brooke (1996)~\cite{sus-brooke1996},
alineado con los conceptos de usabilidad de
ISO~9241-11:2018~\cite{iso9241-11} (eficacia, eficiencia y satisfacci\'on
en un contexto de uso), sin modificar, se aplic\'o a \textbf{18
registros} (P01--P18, ampliado desde los 10 de la Tercera Entrega), con
una tarea de referencia com\'un (login, alta, edici\'on, eliminaci\'on
l\'ogica, logout). Los 18 registros y sus respuestas Q1--Q10 son
evidencia t\'ecnicamente verificable en el repositorio, matem\'aticamente
reproducible (\texttt{scripts/analisis-sus.py}); que corresponden a
participantes reales es una \textbf{declaraci\'on del equipo}. El
protocolo de \texttt{docs/etica/ETHICS.md} exige consentimiento informado
firmado por participante, pero una auditor\'ia documental posterior a
esta entrega constat\'o que esos formularios individuales no est\'an
actualmente disponibles como evidencia verificable; esta situaci\'on se
document\'o mediante una constancia de regularizaci\'on firmada por los
integrantes del proyecto
(\texttt{docs/etica/regularizacion-sus/CONSTANCIA-REGULARIZACION-SUS-BIOPET-2026-08-31.pdf}),
que \textbf{no sustituye} consentimientos individuales.

\begin{table}[H]
  \centering
  \caption{Resultados agregados del instrumento SUS (n = 18).}
  \label{tab:sus-final}
  \begin{tabular}{lr}
    \toprule
    M\'etrica & Valor \\
    \midrule
    Media (SUS Score) & 74{,}44 / 100 \\
    Desviaci\'on t\'ipica (muestral, $n-1$) & 22{,}35 \\
    Intervalo de confianza 95\,\% & $[63{,}33,\ 85{,}56]$ (margen $\pm 11{,}12$) \\
    Mediana (p50) & 82{,}50 \\
    M\'inimo / M\'aximo & 22{,}50 / 97{,}50 \\
    Clasificaci\'on cualitativa (Bangor et al.) & Bueno \\
    \bottomrule
  \end{tabular}
\end{table}

La figura~\ref{fig:sus-resultados-final} presenta la distribuci\'on de
los 18 puntajes SUS reales (diagrama de caja con los puntos individuales
superpuestos), regenerada autom\'aticamente desde \texttt{sus-raw.csv}
mediante el mismo script, sin captura manual.

% EVIDENCIA-JAIME-06 (regenerada automaticamente por Jaime a partir del
% hallazgo del docente sobre esta figura; sustituye, solo en esta
% referencia del informe, a la figura historica
% figuras/zaida/05-sus-resultados.png -- que se conserva SIN modificar en
% el repositorio como artefacto historico de Zaida, no se sobrescribe ni
% se borra)
% Ruta esperada: figuras/jaime/06-sus-resultados-final.png
\evidencia{figuras/jaime/06-sus-resultados-final.png}%
{Distribuci\'on de los 18 puntajes SUS reales: media 74{,}44, DT 22{,}35, mediana 82{,}50, IC 95\,\% $[63{,}33,\ 85{,}56]$ (regenerada por \texttt{docs/informe/scripts/generar-figuras-evidencia.py} desde \texttt{sus-raw.csv}).}%
{fig:sus-resultados-final}

\textbf{Interpretaci\'on.} La media obtenida (74{,}44) se ubica en la
categor\'ia \textbf{Bueno} de la escala de adjetivos SUS, con un IC~95\,\%
que incluye el umbral de referencia de 68 puntos considerado ``sobre el
promedio'' de la industria seg\'un el benchmark emp\'irico de Bangor,
Kortum \& Miller~\cite{bangor2008sus}. El participante con menor puntaje
(P08, 22{,}5) declar\'o experiencia web ``ninguna'', consistente con el
efecto conocido de curva de aprendizaje inicial en usuarios sin experiencia
digital previa. La ampliaci\'on de la muestra de $n=10$ a $n=18$ redujo el
margen del intervalo de confianza (de un margen mayor en la Tercera
Entrega a $\pm 11{,}12$ aqu\'i), pero \textbf{un tama\~no muestral mayor no
implica, por s\'i solo, representatividad externa}: los participantes
siguen siendo reclutados por conveniencia (secci\'on~\ref{cap:amenazas-limitaciones-final}).

\section{Calidad web autom\'atica (Lighthouse)}
\label{sec:lighthouse-final}

La configuraci\'on de auditor\'ia declara umbrales Performance~$\geq$~80,
Accessibility~$\geq$~90, Best Practices~$\geq$~90 y SEO~$\geq$~90~\cite{lighthouse-docs}.
El equipo ejecut\'o tres corridas sucesivas: la inicial del 2026-08-01
(solo perfil m\'ovil, contra \texttt{localhost}, SEO~82/100, umbral no
cumplido); una corrida completa el 2026-08-18
(\texttt{lighthouserc.js}/\texttt{lighthouserc.desktop.js}, tambi\'en
contra \texttt{localhost:4200}, ambos perfiles, SEO~100/100); y la
\textbf{corrida definitiva contra el despliegue p\'ublico de Render}
el \textbf{2026-09-03}
(\texttt{lighthouserc.render.js}/\texttt{lighthouserc.render.desktop.js},
\texttt{docs/mediciones/lighthouse/lhci-20260903-2102-render-*.json},
metadatos en \texttt{lhci-20260903-2102-render.meta.txt}; metodolog\'ia y
medias reproducibles en \texttt{docs/mediciones/lighthouse/RENDER-REPORT.md}),
la primera de las tres auditando \texttt{requestedUrl} p\'ublico
(\texttt{https://biopet-frontend.onrender.com}) en vez de \texttt{localhost}.
Cubre ambos perfiles (m\'ovil y escritorio), 3~corridas por ruta y por
perfil (12~corridas totales) sobre las rutas \texttt{/login} y
\texttt{/mascotas} (ninguna es la portada). El script
\texttt{scripts/run-lighthouse-render.sh} verifica, antes de auditar, que
el despliegue responda \texttt{200} en ambas rutas (sin despliegue en
curso), y valida el campo \texttt{requestedUrl} de cada JSON generado
contra el dominio p\'ublico antes de archivarlo.

\begin{table}[H]
  \centering
  \caption{Resultado real de la corrida Lighthouse del 2026-09-03 contra el despliegue p\'ublico de Render (12 corridas, mobile+desktop, medias calculadas por \texttt{compute-render-averages.mjs}; fuente: \texttt{docs/mediciones/lighthouse/lhci-20260903-2102-render-*.json}).}
  \label{tab:lighthouse-final}
  \begin{tabular}{@{}llrrl@{}}
    \toprule
    Categor\'ia & Perfil & Ruta & Media (0--100) & Umbral \\
    \midrule
    Performance & Escritorio & /login & 99{,}7 & $\geq$ 80 --- Cumplido \\
    Performance & Escritorio & /mascotas & 100{,}0 & $\geq$ 80 --- Cumplido \\
    Performance & M\'ovil & /login & 100{,}0 & $\geq$ 80 --- Cumplido \\
    Performance & M\'ovil & /mascotas & 99{,}0 & $\geq$ 80 --- Cumplido \\
    Accessibility & Ambos & Ambas & 91{,}0 & $\geq$ 90 --- Cumplido \\
    Best Practices & Ambos & /login & 100{,}0 & $\geq$ 90 --- Cumplido \\
    Best Practices & Ambos & /mascotas & 96{,}0 & $\geq$ 90 --- Cumplido \\
    SEO & Ambos & Ambas & \textbf{100{,}0} & $\geq$ 90 --- \textbf{Cumplido} \\
    \bottomrule
  \end{tabular}
\end{table}

% EVIDENCIA-ZAIDA-04
% Ruta esperada: figuras/zaida/04-lighthouse.png
\evidencia{figuras/zaida/04-lighthouse.png}%
{Reporte de Lighthouse CI del frontend de BIOPET contra el despliegue p\'ublico de Render (2026-09-03), por perfil y ruta.}%
{fig:lighthouse-final}

\textbf{Estado real para v1.0.0.} Las \textbf{cuatro} categor\'ias cumplen
su umbral en las 12 corridas del 2026-09-03, en ambos perfiles y ambas
rutas, contra el dominio p\'ublico real (\texttt{requestedUrl} y
\texttt{finalUrl} de los 12 JSON coinciden con
\texttt{https://biopet-frontend.onrender.com/login} o
\texttt{.../mascotas}, sin redirecci\'on). Performance se mantiene en el
rango 99{,}0--100{,}0 en los cuatro grupos perfil/ruta, mejorando el rango
m\'ovil 90--94 de la corrida local del 2026-08-18. Accessibility (91{,}0),
Best Practices (96{,}0--100{,}0) y SEO (100{,}0) se mantienen id\'enticos a
la corrida local previa. Las corridas del 2026-08-01 y 2026-08-18 quedan
como evidencia hist\'orica del proceso de mejora (SEO corregido de 82 a
100) y se conservan sin modificar en
\texttt{docs/mediciones/lighthouse/}; la corrida del 2026-09-03 contra
Render es la que se reporta como definitiva por ser la \'unica auditada
contra el sistema realmente evaluable en producci\'on.

\section{Diferenciaci\'on de los tipos de medici\'on}

Este informe distingue expl\'icitamente: rendimiento (k6, agregado por
corrida, IC~95\,\% con distribuci\'on t), cach\'e (verificaci\'on aislada
por clave con \texttt{redis-cli}), cobertura (JaCoCo, LINE/BRANCH/COMPLEXITY
sobre el \texttt{BUNDLE} completo), usabilidad (SUS, ejecutada, $n=18$),
calidad web (Lighthouse, ejecutada en ambos perfiles mobile y desktop, 12
corridas, cuatro categor\'ias cumplidas), y seguridad (revisi\'on manual OWASP +
an\'alisis din\'amico ZAP + an\'alisis est\'atico SpotBugs/Find Security
Bugs, cap\'itulo~\ref{cap:seguridad}), para evitar mezclar cifras de
naturaleza distinta.
